\chapter*{Introducción General}
\phantomsection
\addcontentsline{toc}{chapter}{Introducción General}

El presente documento consolida el trabajo realizado durante las Actividades de Aprendizaje 4 y 5 de la Fase de Análisis, constituyendo el paquete de \textbf{Especificación y Validación de Requisitos} para el proyecto ``Gemelo Digital de Mantenimiento''.

A través de las cuatro evidencias aquí presentadas, se sigue un flujo de trabajo que va desde la abstracción de alto nivel hasta la validación visual y funcional. Se inicia con la \textbf{traducción de las necesidades del negocio en requisitos técnicos formales (AA4-EV01)}, se profundiza en la \textbf{especificación detallada mediante el estándar IEEE 830 y Historias de Usuario (AA4-EV02)}, se documenta el \textbf{entorno de gestión de dichos requisitos (AA5-EV01)} y, finalmente, se \textbf{validan los conceptos clave mediante prototipos y casos de prueba (AA5-EV02)}.

El objetivo final de este documento es entregar una ``fuente de verdad'' técnica y unificada que servirá como base sólida para la fase de Diseño y, posteriormente, la Ejecución del software.
